% Kandydat restrukturyzacji; nie jest aktywnym rozdziałem Atlasu.
\section{Splot funkcji i nierówność Younga}
\providecommand{\atlasref}[2]{#2}

\subsection{Nierówność Younga dla splotu}

\begin{twierdzenie}[Nierówność Younga]
\label{thm:dt-young-general}
Niech $1\le p,q,r\le\infty$ i
$1+1/r=1/p+1/q$, z konwencją $1/\infty=0$.
Dla $f\in L^p(\R^d)$ i $g\in L^q(\R^d)$ całka
\[
 (f*g)(x)=\int_{\R^d}f(y)g(x-y)\,\mathrm d\lambda_d(y)
\]
jest bezwzględnie zbieżna dla prawie każdego $x$,
określa element $L^r$ niezależny od reprezentantów
i spełnia $\|f*g\|_r\le\|f\|_p\|g\|_q$.
\end{twierdzenie}
\begin{proof}
Dla $r=\infty$ mamy $1/p+1/q=1$ i bezpośrednio
stosujemy nierówność Höldera do całki po $y$.
Niech zatem $r<\infty$. Wtedy $p,q<\infty$ oraz
$p,q\le r$. Rozkładamy moduł iloczynu na trzy czynniki:
\[
 |f(y)g(x-y)|
 =\bigl(|f(y)|^p|g(x-y)|^q\bigr)^{1/r}
 |f(y)|^{1-p/r}|g(x-y)|^{1-q/r}.
\]
Jeżeli któryś z dwóch ostatnich wykładników jest zerowy,
odpowiedni czynnik zastępujemy jednością.
Stosujemy nierówność Höldera dla trzech czynników,
z odwrotnościami wykładników
$1/r$, $1/p-1/r$ oraz $1/q-1/r$.
Ich suma jest równa jedności z założenia.
Otrzymujemy
\[
 \int_{\R^d}|f(y)g(x-y)|\,\mathrm d\lambda_d(y)
 \le A(x)^{1/r}\|f\|_p^{1-p/r}\|g\|_q^{1-q/r},
\]
gdzie
$A(x)=\int_{\R^d}|f(y)|^p|g(x-y)|^q\,\mathrm d\lambda_d(y)$.
Z twierdzenia Tonellego
$\int_{\R^d} A(x)\,\mathrm d\lambda_d(x)=\|f\|_p^p\|g\|_q^q$,
więc $A(x)<\infty$ prawie wszędzie.
Podnosimy oszacowanie do potęgi $r$, całkujemy po $x$
i otrzymujemy $\|f*g\|_r^r\le\|f\|_p^r\|g\|_q^r$.
Przypadek zerowej normy daje splot zerowy prawie wszędzie.
Mierzalność i niezależność od reprezentantów wynikają
z twierdzenia Tonellego zastosowanego do modułów i
do zbiorów, na których zmieniono reprezentanty.
\end{proof}

\begin{uwaga}[Dwie nierówności Younga i ostrość]
\label{rem:dt-young-scope}
Nie należy utożsamiać nierówności Younga dla splotu
z nierównością dla liczb dodatnich
$ab\le a^p/p+b^{p'}/p'$.
Stała jeden w Twierdzeniu~\ref{thm:dt-young-general}
wystarcza do konstrukcji ciągłych operatorów splotowych;
w ogólności nie jest ona optymalna. Przypadek $p=1$
obejmuje oszacowanie stosowane przy przybliżaniu jedynki.
\end{uwaga}
